\begin{table}[h]
    \caption{Finance terms used in this paper.}
    \label{tab:glossary}
    \centering
    \small
    \begin{tabular}{@{}p{0.22\linewidth} p{0.45\linewidth} p{0.27\linewidth}@{}}
      \toprule
      Term & Meaning & Answers the question \\
      \midrule
      Spot price $S$
        & The underlying's price today. We price $C/S$, so its value never matters.
        & What does the stock cost now? \\
      Call option
        & A contract giving the right, but not the obligation, to buy the underlying at a fixed price $K$ at a future time $T$.
        & What is being priced? \\
      Call price $C$
        & What the call option costs today. The target $f(x)$ the transformer predicts (as $C/S$, standardised).
        & How much is the right to buy worth? \\
      Strike $K$
        & The fixed price at which the call lets you buy.
        & At what price can I buy? \\
      Log-moneyness $m$
        & $m = \ln(K/S)$: the strike relative to today's price. $m = 0$ means $K = S$.
        & How far is the strike from today's price? \\
      Maturity $T$
        & Time until the option expires, in years.
        & How long until the option ends? \\
      Risk-free rate $r$
        & Interest earned on a riskless investment; fixed at $r = 0.02$.
        & What would money earn risk-free meanwhile? \\
      Volatility $\sigma$
        & How much the underlying's price fluctuates (annualised std of log-returns). The hidden parameter per prompt.
        & How uncertain is the future price? \\
      Black--Scholes formula
        & The standard formula giving $C$ from $(S, K, T, r, \sigma)$. Only $\sigma$ is not directly observable.
        & What is the fair price of the option? \\
      Implied volatility
        & The $\sigma$ that makes the Black--Scholes price equal an observed price.
        & What volatility is the market price assuming? \\
      Volatility smile / skew
        & Implied volatility varying with the strike instead of being constant; equity markets show higher $\sigma$ for low strikes (skew).
        & Does one $\sigma$ fit all strikes? \\
      Calibration
        & Fitting the volatility parameters so model prices match observed prices, then using them to price new options.
        & Which $\sigma$ explains the prices I see? \\
      Bid--ask spread
        & The gap between buying and selling prices in a market; makes observed prices noisy. Motivates our noisy evaluation.
        & How exact are observed prices? \\
      \bottomrule
    \end{tabular}
  \end{table}